% Research section prepared by harmonic_jets agent.
% Requires amsmath, amssymb, amsthm; theorem/corollary/proposition/remark envs.
% Citation keys proposed in harmonic_sources.bib (or adapt to main bibliography).
\section{Fully shifted height-one sums: an entire difference and exact Laurent data}
\label{xid:hsh:sec:main}

The inner and outer shifts of a nested harmonic sum are different pieces
of data. The earlier shifted height-one families in this project move the
outer denominator while retaining ordinary inner harmonic numbers. The
newest nested-jet report instead gives all nested denominators a common
variable exponent. Here every denominator has the same shift, while only
the outer exponent varies. A Gauss connection formula shows that changing
the common shift preserves a universal singular part, up to an explicit
Gamma factor. The remaining function is entire in the spectral variable,
and its values at every nonpositive integer are finite Stirling transforms.

At depth one this family belongs to the harmonic Hurwitz-zeta literature
of Karg\i n, Dil, Cenkci and Can \cite{hsh:Kargin2023}; the recent preprint
of Kara \"Ozt\"urk and Can treats the two-parameter depth-one family and
negative Laurent coefficients \cite{hsh:KaraCan2026}. The latter comparison
is based on its retrieved abstract: its complete text was not available
in this research environment. Thus the depth-one meromorphic continuation
and negative-integer evaluation problem are not claimed as new. The
all-depth entire difference, finite Stirling rule, and complete pole-order
classification below were derived here and were not found in the
inspected repository material; priority over all literature is not asserted.

\subsection{The common-shift convention}

Let $\Re a>0$. All powers of $n+a$ use the principal logarithm, and $u$
will remain in a sufficiently small disc centered at zero. Set
\begin{align}
 e_r(n;a)&=[u^r]\prod_{j=0}^{n-1}\left(1+\frac{u}{j+a}\right),
       &e_0(n;a)&=1,\label{xid:hsh:eq:er}\\
 M_r(s;a)&=\sum_{n=0}^{\infty}\frac{e_r(n;a)}{(n+a)^s},
       &&\Re s>1,\label{xid:hsh:eq:Mr}\\
 F_a(s,u)&=\sum_{n=0}^{\infty}
       \frac{(a+u)_n}{(a)_n(n+a)^s}
       =\sum_{r\ge0}M_r(s;a)u^r.
       \label{xid:hsh:eq:F}
\end{align}
The first product is empty when $n=0$. Equivalently,
\[
 M_r(s;a)=\sum_{n_0>n_1>\cdots>n_r\ge0}
       \frac{1}{(n_0+a)^s(n_1+a)\cdots(n_r+a)}.
\]
Thus $M_0(s;a)=\zeta(s,a)$ and $M_r(s;1)=\zeta(s,\{1\}^r)$.
If $H_n^{(k)}(a)=\sum_{j=0}^{n-1}(j+a)^{-k}$, then
\[
 e_1(n;a)=H_n^{(1)}(a),\qquad
 e_2(n;a)=\frac{(H_n^{(1)}(a))^2-H_n^{(2)}(a)}2.
\]
The general $e_r$ is the corresponding elementary-symmetric Bell
polynomial. Uniformly on compact $a$-sets,
$e_r(n;a)=O((1+\log n)^r)$; hence every fixed-depth series converges
normally for $\Re s>1$. The Gamma-ratio Dirichlet series in \eqref{xid:hsh:eq:F} converges normally
when $\Re s>1+\Re u$, locally uniformly in its parameters. Interchange
with its depth Taylor series is justified, for example, on
$\Re s>1+|u|$; at each fixed $\Re s>1$ one may choose a sufficiently
small disc in $u$. The direct-series estimate follows from
\[
 \frac{(a+u)_n}{(a)_n}
 =\frac{\Gamma(a)}{\Gamma(a+u)}
    \frac{\Gamma(n+a+u)}{\Gamma(n+a)}
 =O(n^{\Re u}).
\]
Gamma-ratio asymptotics and their uniform sector versions are classical
\cite{hsh:DLMFGamma}.

Introduce
\begin{align}
 C_a(u)&=\frac{\Gamma(a)}{\Gamma(a+u)}
        =\sum_{j\ge0}c_j(a)u^j,\label{xid:hsh:eq:Ca}\\
 K_a(u)&=\Gamma(1+u)C_a(u)
        =\sum_{j\ge0}k_j(a)u^j.\label{xid:hsh:eq:Ka}
\end{align}
The Gamma logarithm is analytic on the right half-plane, so
\begin{equation}\label{xid:hsh:eq:Cexp}
 C_a(u)=\exp\left(-\sum_{j\ge1}
                \frac{\psi^{(j-1)}(a)}{j!}u^j\right).
\end{equation}
Consequently every coefficient is a finite polynomial in polygamma
functions. In particular, with $P=\psi(a)$, $Q=\psi'(a)$, and
$R=\psi''(a)$,
\[
 c_0=1,\quad c_1=-P,\quad c_2=\frac{P^2-Q}{2},\quad
 c_3=\frac{-P^3+3PQ-R}{6}.
\]
The standard polygamma--Hurwitz identities can convert these coordinates
to zeta values when useful \cite{hsh:DLMFPolygamma,DLMFhurwitz}.

\subsection{The exact entire difference}

\begin{theorem}[Universal singular part under a common shift]
\label{xid:hsh:thm:entire}
For $\Re a>0$ and $u$ in a sufficiently small disc centered at zero,
\begin{equation}\label{xid:hsh:eq:D}
 D(s,u;a):=F_a(s,u)-K_a(u)F_1(s,u)
\end{equation}
extends jointly holomorphically to all $s\in\mathbb C$, locally in $a,u$.
Consequently, for every $r\ge0$,
\begin{equation}\label{xid:hsh:eq:Dr}
 D_r(s;a):=M_r(s;a)-\sum_{j=0}^rk_j(a)M_{r-j}(s;1)
\end{equation}
is entire in $s$.

For every integer $m\ge0$, its complete depth-generating value is
\begin{equation}\label{xid:hsh:eq:T}
 \boxed{\quad D(-m,u;a)
 =-\sum_{j=1}^{m+1}
       {\left\{\begin{matrix}m+1\\j\end{matrix}\right\}}
       \frac{(a-1)^{\underline j}}{u+j}.\quad}
\end{equation}
Here braces denote Stirling numbers of the second kind and
$x^{\underline j}=x(x-1)\cdots(x-j+1)$.
Equivalently,
\begin{equation}\label{xid:hsh:eq:Delta}
 \boxed{\quad
 D_r(-m;a)=\Delta_{m,r}(a):=
 (-1)^{r+1}\sum_{j=1}^{m+1}
       {\left\{\begin{matrix}m+1\\j\end{matrix}\right\}}
       \frac{(a-1)^{\underline j}}{j^{r+1}}.
 \quad}
\end{equation}
Thus every value in \eqref{xid:hsh:eq:Delta} is a rational polynomial in $a$
of degree at most $m+1$.
\end{theorem}

\begin{proof}
In the common half-plane of convergence, the Mellin kernel of $F_a$ is
\begin{equation}\label{xid:hsh:eq:Mellin}
 \Gamma(s)F_a(s,u)=\int_0^{\infty}t^{s-1}G_a(t,u)\,dt,
 \qquad G_a(t,u)=e^{-at}
          {}_2F_1(1,a+u;a;e^{-t}).
\end{equation}
Termwise integration follows from absolute convergence. The Gauss
connection formula at $z=1$ \cite{hsh:DLMFGauss} simplifies here to
\begin{equation}\label{xid:hsh:eq:kernelconnection}
 \begin{split}
 G_a(t,u)&=K_a(u)\frac{e^{-t}}{(1-e^{-t})^{1+u}}+R_a(t,u),\\
 R_a(t,u)&=-\frac{a-1}{1+u}e^{-at}
       {}_2F_1(1,a+u;2+u;1-e^{-t}).
 \end{split}
\end{equation}
Indeed its second local hypergeometric factor is
${}_2F_1(a-1,-u;-u;1-z)=z^{1-a}$, and its two connection
coefficients reduce to $K_a(u)$ and $-(a-1)/(1+u)$.
The identity is first read for noninteger $u$, where the generic
connection formula applies, and then continued through $u=0$.
The simplified two terms are individually analytic in $u$ near zero.

The function $R_a(t,u)$ is analytic in $t$ near zero, jointly in the
other parameters. Both $G_a$ and the first term in
\eqref{xid:hsh:eq:kernelconnection} decay exponentially as $t\to+\infty$,
locally uniformly for $\Re a>0$. Thus the same holds for $R_a$.
Write $R_a(t,u)=\sum_{j\ge0}r_j(a,u)t^j$. For any $L\ge1$,
\[
 \begin{split}
 \Gamma(s)D(s,u;a)={}&\sum_{j=0}^{L-1}\frac{r_j(a,u)}{s+j}\\
 &+\int_0^1t^{s-1}\left(R_a(t,u)-\sum_{j=0}^{L-1}r_j(a,u)t^j\right)dt
  +\int_1^{\infty}t^{s-1}R_a(t,u)\,dt.
 \end{split}
\]
The remainder integrals are jointly holomorphic for $\Re s>-L$.
All their parameter derivatives have the same local integrability,
possibly with fixed powers of $\log t$. Each displayed simple pole
at $s=-j$ is canceled by the simple zero of $1/\Gamma(s)$ there.
Increasing $L$ proves joint entireness in $s$ and gives
\begin{equation}\label{xid:hsh:eq:Dlocal}
 D(-m,u;a)=(-1)^m m!\,[t^m]R_a(t,u).
\end{equation}
For fixed $m$, the right side is a polynomial in $a$ of degree at most
$m+1$: only the first $m+1$ terms of the local hypergeometric series
can contribute, and the $j$-th term has a factor $(1-e^{-t})^j$.
Its coefficients are rational functions of $u$, analytic near zero.

It remains to identify this polynomial. For a positive integer $A$,
shifting the index of the defining series gives the exact finite
identity
\[
 D(s,u;A)=-K_A(u)\sum_{k=1}^{A-1}
           \frac{(1+u)_{k-1}}{(k-1)!\,k^s}.
\]
At $s=-m$, use
\[
 k^m=\sum_{j=1}^{m+1}
       {\left\{\begin{matrix}m+1\\j\end{matrix}\right\}}
         (k-1)^{\underline{j-1}}.
\]
The finite binomial summation yields
\[
 K_A(u)\sum_{k=1}^{A-1}
       \frac{(1+u)_{k-1}}{(k-1)!}
              (k-1)^{\underline{j-1}}
       =\frac{(A-1)^{\underline j}}{u+j}.
\]
For completeness, after putting $\ell=k-j$, the nonzero summands are
$(1+u)_{j-1}(u+j)_\ell/\ell!$; their sum is
$(1+u)_{j-1}(u+j+1)_{A-j-1}/(A-j-1)!$.
Multiplication by $K_A=(A-1)!/(1+u)_{A-1}$ gives the stated quotient.
When $A\le j$, both sides are zero. Therefore \eqref{xid:hsh:eq:T} holds
at infinitely many positive integers $A$. Both sides are polynomials
in $a$ of bounded degree, so they agree identically. Finally expand
$(u+j)^{-1}=\sum_{r\ge0}(-1)^ru^r/j^{r+1}$ to obtain
\eqref{xid:hsh:eq:Delta}.
\end{proof}

At $r=0$, the theorem reduces to the usual Bernoulli-polynomial
identity for $\zeta(-m,a)-\zeta(-m)$. At positive depths it gives a
finite rational interpolation of the entire difference after the
Gamma mixing of depths. The theorem does not identify
$M_r(-m;a)$ with a convergent series: the defining series diverges
there, and the statement concerns its meromorphic continuation.

\subsection{A bridge to independent inner and outer shifts}

The same argument unifies the earlier outer-shift calculus with the
common-shift family. Let $\Re a>0$, $\Re b>0$, and
$\Re c>0$, where $c=1+b-a$. Define
\[
 F_{a,b}(s,u)=\sum_{n\ge0}\frac{(a+u)_n}{(a)_n(n+b)^s},
 \qquad B_c(s,u)=\sum_{n\ge0}\frac{(1+u)_n}{n!(n+c)^s}.
\]
The function $B_c$ is precisely the previous generator with an outer
shift only. At $a=1$, $F_{a,b}=B_b$; at $b=a$, $F_{a,b}=F_a$.

\begin{corollary}[Two shifts and a finite binomial correction]
\label{xid:hsh:cor:twoshifts}
The difference
\[
 D_{a,b}(s,u)=F_{a,b}(s,u)-K_a(u)B_{1+b-a}(s,u)
\]
is entire in $s$, locally in its other parameters. For $m\ge0$,
\begin{equation}\label{xid:hsh:eq:twoshifts}
 D_{a,b}(-m,u)
   =\sum_{v=0}^m\binom mv(b-a)^{m-v}D(-v,u;a).
\end{equation}
Every term on the right is the explicit finite Stirling sum in
\eqref{xid:hsh:eq:T}. In depth $r$ the corresponding correction is
\[
 [u^r]D_{a,b}(-m,u)
    =\sum_{v=0}^m\binom mv(b-a)^{m-v}\Delta_{v,r}(a).
\]
\end{corollary}
\begin{proof}
Multiplication of the Mellin kernel in \eqref{xid:hsh:eq:kernelconnection}
by $e^{-(b-a)t}$ gives
\[
 e^{-bt}{}_2F_1(1,a+u;a;e^{-t})
   =K_a(u)\frac{e^{-(1+b-a)t}}{(1-e^{-t})^{1+u}}
         +e^{-(b-a)t}R_a(t,u).
\]
The last kernel is analytic at zero and decays exponentially at infinity
under the stated hypotheses. The same Mellin subtraction proves
entireness. Taking its $m$-th Taylor coefficient and applying
\eqref{xid:hsh:eq:Dlocal} proves \eqref{xid:hsh:eq:twoshifts} by the ordinary
binomial theorem.
\end{proof}

Thus no new singular function is required to separate arbitrary inner
and outer shifts in this domain: the existing outer-shift generator
carries the singular terms and a fully explicit polynomial correction
accounts for the values of the entire remainder at nonpositive integers.
The higher spectral derivatives of this remainder remain a separate
identity problem.

\subsection{The complete Laurent principal part and finite part}

Define the rational polynomials $b_j(u)$ and $W_m(u)$ by
\begin{align}
 e^{-t}\left(\frac{t}{1-e^{-t}}\right)^{1+u}
       &=\sum_{j\ge0}b_j(u)t^j,\label{xid:hsh:eq:b}\\
 W_m(u)&=(u-m)_{m+1}b_{m+1}(u).
       \label{xid:hsh:eq:W}
\end{align}
The logarithm of the parenthesized factor is the analytic germ at zero
whose value there is zero. Each $W_m$ is a rational polynomial divisible
by $u$. The first three are
\begin{align}
 W_0(u)&=\frac{u(u-1)}2,\label{xid:hsh:eq:W0}\\
 W_1(u)&=\frac{u(u-1)(u-2)(3u-1)}{24},\\
 W_2(u)&=\frac{u^2(u-1)^2(u-2)(u-3)}{48}.
\end{align}

\begin{theorem}[All depths, every nonpositive integer]
\label{xid:hsh:thm:Laurent}
Let $m,r\ge0$ and $\Re a>0$. Then
\begin{equation}\label{xid:hsh:eq:Laurent}
 \boxed{\quad
 M_r(-m+\varepsilon;a)=
 \sum_{h=0}^r
     \frac{[u^{r+1-h}]\{C_a(u)W_m(u)\}}{\varepsilon^h}
   +\Delta_{m,r}(a)+O(\varepsilon).
 \quad}
\end{equation}
In particular,
\begin{equation}\label{xid:hsh:eq:FP}
 \operatorname{FP}_{s=-m}M_r(s;a)
  =[u^{r+1}]\{C_a(u)W_m(u)\}+\Delta_{m,r}(a).
\end{equation}
The finite part belongs to
$\mathbb Q[a,\psi(a),\psi'(a),\ldots,\psi^{(r-1)}(a)]$,
where at $r=0$ the polygamma list is empty.
\end{theorem}

\begin{proof}
For the universal unshifted kernel, Mellin subtraction gives, near
$(s,u)=(-m,0)$,
\[
 F_1(-m+\varepsilon,u)
 =\frac{b_{m+1}(u)}{\Gamma(-m+\varepsilon)(\varepsilon-u)}
       +\frac{Q_m(\varepsilon,u)}{\Gamma(-m+\varepsilon)},
\]
where $Q_m$ is jointly holomorphic. The last term contributes only
$O(\varepsilon)$ after taking any fixed coefficient in $u$.
Expand $1/\Gamma(-m+\varepsilon)$ in powers of $\varepsilon$;
its constant coefficient is zero. From the first term, the coefficient
of $u^r\varepsilon^{-h}$, for $h\ge0$, is
\[
 [u^{r+1-h}]\frac{K_a(u)b_{m+1}(u)}{\Gamma(u-m)}
   =[u^{r+1-h}]\{C_a(u)W_m(u)\}.
\]
The second equality is the Gamma recurrence. Finally add the entire
function $D(s,u;a)$ and apply \eqref{xid:hsh:eq:Delta}. Since $W_m$ has
no constant coefficient, at most the first $r$ derivatives of
$\log\Gamma(a+u)$ can occur, giving the stated algebra.
\end{proof}

The polynomial correction in \eqref{xid:hsh:eq:FP} is essential. For instance,
with $r=m=0$, omitting it gives $-1/2$, whereas the correct value is
$\zeta(0,a)=1/2-a$. Thus the singular-coefficient rule for the earlier
outer-shift family cannot be transferred to the fully shifted family
without its analytic-kernel contribution. This is an extension guard,
not a claim that the earlier report made this incorrect transfer.

\begin{corollary}[Exact pole orders]
\label{xid:hsh:cor:poles}
At $s=1$, the order of the pole of $M_r$ is $r+1$.
At $s=0$ and every negative odd integer $s=-m$, the order is exactly $r$
for $r\ge1$, with leading coefficient $\zeta(-m)$.
At a negative even integer $s=-m$, $m\ge2$, the order is exactly $r-1$
for $r\ge2$, with leading coefficient
\begin{equation}\label{xid:hsh:eq:evenleading}
 \frac{m+1}{2m}B_m.
\end{equation}
At these negative even integers, $M_1$ is holomorphic.
All the specified leading coefficients are independent of $a$.
\end{corollary}

\begin{proof}
The assertion at one follows from the next theorem. From
\eqref{xid:hsh:eq:b},
\[
 [u]W_m(u)=(-1)^m m!\,b_{m+1}(0)=\zeta(-m).
\]
This is nonzero for $m=0$ and for positive odd $m$. If $m\ge2$ is even,
$b_{m+1}(0)=B_{m+1}/(m+1)!=0$. The two elementary expansions
\[
 \frac{t}{e^t-1}=1-\frac t2+
       \sum_{q\ge1}\frac{B_{2q}}{(2q)!}t^{2q},\qquad
 \log\frac{t}{1-e^{-t}}
       =\frac t2-\sum_{q\ge1}
          \frac{B_{2q}}{2q(2q)!}t^{2q}
\]
give
\[
 [t^{m+1}]\frac{t}{e^t-1}\log\frac{t}{1-e^{-t}}
       =\frac{m+1}{2m}\frac{B_m}{m!}.
\]
Multiplication by the leading term $m!u$ of $(u-m)_{m+1}$ proves
$[u^2]W_m=(m+1)B_m/(2m)$. This is nonzero, since every positive even
Bernoulli number is nonzero \cite{hsh:DLMFBernoulli}. Thus the zero of
$W_m$ at zero is exactly double in this case and simple in the preceding
case. Since $C_a(0)=1$ and $D$ is entire, the pole orders and leading
coefficients follow directly from \eqref{xid:hsh:eq:Laurent}.
\end{proof}

\subsection{The principal Stieltjes point}

\begin{theorem}[Gamma coefficients at the principal pole]
\label{xid:hsh:thm:one}
For $r\ge0$,
\begin{equation}\label{xid:hsh:eq:one}
 M_r(1+\varepsilon;a)
 =\sum_{j=0}^r\frac{c_j(a)}{\varepsilon^{r+1-j}}
       +c_{r+1}(a)+O(\varepsilon).
\end{equation}
Also
\begin{equation}\label{xid:hsh:eq:D1}
 D(1,u;a)=\frac{K_a(u)-1}{u},\qquad D_r(1;a)=k_{r+1}(a).
\end{equation}
The apparent singularity on the right is removable.
\end{theorem}

\begin{proof}
Let $p_n=(a+u)_n/(a)_n$. Its exact difference is
$p_{n+1}-p_n=up_n/(n+a)$. For small $\Re u<0$, $p_n\to0$;
therefore telescoping proves $F_a(1,u)=-1/u$.
The leading kernel term in \eqref{xid:hsh:eq:kernelconnection} shows that
\[
 F_a(1+\varepsilon,u)=\frac{C_a(u)}{\varepsilon-u}
                         +Q_a(\varepsilon,u)
\]
with $Q_a$ jointly holomorphic near zero. Evaluate at $\varepsilon=0$
for $\Re u<0$ to obtain $Q_a(0,u)=(C_a(u)-1)/u$, and continue
holomorphically through $u=0$. Coefficient extraction proves
\eqref{xid:hsh:eq:one}. Applying the same telescoping value to the
entire difference gives \eqref{xid:hsh:eq:D1}.
\end{proof}

At depth one, the commonly used inclusive harmonic Hurwitz function is
\[
 \zeta_H(s,a)=\sum_{n\ge0}\frac{\sum_{j=0}^{n}(j+a)^{-1}}{(n+a)^s}
            =M_1(s;a)+\zeta(s+1,a).
\]
Thus \eqref{xid:hsh:eq:one} gives
\[
 \zeta_H(1+\varepsilon,a)
 =\frac1{\varepsilon^2}-\frac{\psi(a)}{\varepsilon}
       +\frac{\psi(a)^2+\psi'(a)}2+O(\varepsilon).
\]
The strict and inclusive conventions differ even at the finite part.
The formula at all depths gives the Gamma-polygamma coefficients of the
strict height-one family; higher regular Laurent coefficients are not
claimed to reduce to this same finite algebra.

\subsection{Explicit values and Gamma antiderivatives}

Writing $P=\psi(a)$ and $Q=\psi'(a)$, the first examples are
\begin{align}
 M_1(\varepsilon;a)
  &=-\frac1{2\varepsilon}+a-\frac12+\frac P2+O(\varepsilon),
       \label{xid:hsh:eq:M10}\\
 M_1(-1+\varepsilon;a)
  &=-\frac1{12\varepsilon}
       +\frac{a^2+a}{4}-\frac18+\frac P{12}+O(\varepsilon),\\
 M_1(-2;a)&=\frac{8a^3+6a^2-2a-3}{72},
       \label{xid:hsh:eq:M12}\\
 M_2(\varepsilon;a)
  &=-\frac1{2\varepsilon^2}+\frac{1+P}{2\varepsilon}
       +1-a-\frac P2-\frac{P^2-Q}{4}+O(\varepsilon),\\
 M_2(-1+\varepsilon;a)
  &=-\frac1{12\varepsilon^2}
       +\frac{3/8+P/12}{\varepsilon}
       +\frac{-3a^2-15a+8-9P-P^2+Q}{24}
       +O(\varepsilon).\label{xid:hsh:eq:M21}
\end{align}
At the half shift, put $L=\gamma+2\log2$, so
$\psi(1/2)=-L$ and $\psi'(1/2)=\pi^2/2$. Then
\begin{align}
 M_1(-2;1/2)&=-\frac1{48},\\
 M_2(-2+\varepsilon;1/2)
       &=\frac1{8\varepsilon}+\frac L8-\frac{19}{288}
                    +O(\varepsilon).
\end{align}
In the convergent half-plane, for example,
$M_2(s;1/2)$ is one half the Dirichlet series whose numerator is
$(\sum_{j<n}(j+1/2)^{-1})^2-\sum_{j<n}(j+1/2)^{-2}$.
Its displayed pole and finite part follow without a numerical relation search.

There is an exact antiderivative for every depth-one finite part. Put
\[
 P_m(a)=[u^2]W_m(u)+
       \sum_{j=1}^{m+1}{\left\{\begin{matrix}m+1\\j\end{matrix}\right\}}
                 \frac{(a-1)^{\underline j}}{j^2}\in\mathbb Q[a].
\]
Then
\begin{equation}\label{xid:hsh:eq:depthonegeneral}
 \operatorname{FP}_{s=-m}M_1(s;a)
       =P_m(a)-\zeta(-m)\psi(a),
\end{equation}
and on the right half-plane
\begin{equation}\label{xid:hsh:eq:primitive}
 \int\operatorname{FP}_{s=-m}M_1(s;a)\,da
       =\mathcal P_m(a)-\zeta(-m)\log\Gamma(a)+\text{constant},
 \qquad \mathcal P_m'=P_m.
\end{equation}
For even $m\ge2$, the Gamma term vanishes and the value in
\eqref{xid:hsh:eq:depthonegeneral} is itself a polynomial. For $m=0,1$,
\begin{align*}
 \int\operatorname{FP}_{s=0}M_1(s;a)\,da
       &=\frac{a^2-a}{2}+\frac12\log\Gamma(a)+\text{constant},\\
 \int\operatorname{FP}_{s=-1}M_1(s;a)\,da
       &=\frac{a^3}{12}+\frac{a^2}{8}-\frac a8
                      +\frac1{12}\log\Gamma(a)+\text{constant}.
\end{align*}
Higher derivatives in $a$ consequently give explicit polygamma formulas.
This assertion is about finite Laurent coefficients; it does not replace
a divergent defining series by termwise integration at a negative order.

\subsection{Reproduction and further questions}

The accompanying verification script performs $72$ floating-point checks
at $75$ decimal working digits. Its maximum absolute residual is below
$8.5\times10^{-52}$. Thirty checks compare \eqref{xid:hsh:eq:T} with an
independent continuation of the original Gamma-ratio Dirichlet series,
using its Bernoulli-polynomial asymptotic coefficients and Hurwitz-zeta
tails. Two repeats with a larger cutoff and more tail terms have residuals
below $1.7\times10^{-64}$. The remaining checks audit the Gauss connection,
its mixed Taylor coefficients, and \eqref{xid:hsh:eq:D1}. These are
floating-point checks, not interval certificates; the all-parameter
proof is the analytic argument above.

A constructive next step is to determine the higher regular Taylor
coefficients of the entire difference $D_r(s;a)$ at $s=-m$. Unlike its
value, its first derivative includes a genuine integral of the analytic
kernel and need not be a rational polynomial. A second direction is to
retain independent perturbations of the inner exponents: the finite-product
model then requires more than one spectral direction, and the singular
part need not be governed by a single Gamma quotient. A third direction is to compare the finite
Stirling transforms in \eqref{xid:hsh:eq:Delta} with established
poly-Bernoulli and generalized Bernoulli polynomial systems, with exact
normalization matching before claiming an identification. Finally,
character-weighted sums in the common shift may yield further finite
cyclotomic reductions of the polygamma algebra in \eqref{xid:hsh:eq:FP}.
